% ============================================================
% AGF-ADNet 模型与算法描述
% 说明：主文件（独立编译版本），使用 XeLaTeX 编译。
% ============================================================

\documentclass[UTF8]{ctexart}

\usepackage{amsmath}
\usepackage{amssymb}
\usepackage{xeCJK}

\begin{document}

\section{AGF-ADNet 模型与算法}

\subsection{方法概述}

工业生产过程中通常会安装大量传感器来实时监测系统状态，这些传感器每时每刻都在产生数据，形成了一组\textbf{多变量时间序列}（Multivariate Time Series）。一旦系统出现故障或异常，传感器读数会偏离正常范围。异常检测的目标就是及时发现这种偏离。

本研究提出一种面向工业过程异常检测的自适应图--频率融合异常检测网络（Adaptive Graph-Frequency Anomaly Detection Network，AGF-ADNet）。该模型的核心思想非常直观：

\begin{quote}
\textbf{核心思想}：如果我们训练一个模型，让它学会重构正常状态下的信号，那么当输入异常信号时，模型就无法准确重构，导致\textbf{重构误差}（Reconstruction Error）变大。重构误差越大，输入就越可能是异常。
\end{quote}

这一思路在机器学习中被称为\textbf{基于重构的异常检测}（Reconstruction-based Anomaly Detection），与自编码器（Autoencoder）的基本原理一致：模型只见过正常数据，因此对异常数据的重构能力较差。

AGF-ADNet 包含以下关键组件：
\begin{enumerate}
    \item \textbf{自适应图学习模块（GraphLearner）}：自动学习变量间的依赖关系，生成稀疏邻接矩阵；
    \item \textbf{时域分支}：由图卷积网络（GCN）和多尺度时序卷积网络（Multi-Scale TCN）组成，分别捕获变量间空间依赖和多尺度时序模式；
    \item \textbf{频域分支（FrequencyImputer）}：在频域中通过注意力机制增强信号的频率特征；
    \item \textbf{门控融合模块（GatedFusion）}：自适应地融合时域与频域分支的输出；
    \item \textbf{Transformer 重构器}：基于自注意力机制对融合后的表示进行全局时序建模与信号重构。
\end{enumerate}

模型采用\textbf{自监督学习}（Self-Supervised Learning）范式训练。这意味着我们不需要人工标注哪些样本是异常的——模型通过一个精心设计的\textbf{掩码--重构}任务来自我学习：在训练阶段对已观测数据随机遮盖一部分，迫使模型学习恢复被遮盖的信号，从而建立对正常模式的内在理解。

% ============================================================
\subsection{模型架构}

\subsubsection{张量表示与符号约定}

在深入各模块之前，我们先明确本文使用的数学符号。在深度学习中，数据通常以\textbf{张量}（Tensor）的形式表示——张量可以理解为多维数组的推广。

设输入多变量时间序列为 $\mathbf{X} \in \mathbb{R}^{B \times S \times N}$，其中：
\begin{itemize}
    \item $B$：批量大小（Batch Size），即同时处理的样本数量，主要用于加速训练；
    \item $S$：时间窗口长度（Sequence Length），默认 $S=60$，表示每个样本包含 60 个时间步；
    \item $N$：变量个数（Number of Variables），默认 $N=41$，对应 41 个传感器。
\end{itemize}

直观地说，$\mathbf{X}$ 是一个三维数组，$X_{b,s,n}$ 表示第 $b$ 个样本在第 $s$ 个时间步处第 $n$ 个传感器的读数。

对应的观测掩码为 $\mathbf{M} \in \{0,1\}^{B \times S \times N}$，其中 $M_{b,s,n}=1$ 表示该位置的值已被观测（数据可用），$M_{b,s,n}=0$ 表示缺失。掩码的作用类似于一张"过滤网"：我们只在有真实数据的位置上计算损失和评估模型。

\subsubsection{前向传播总览}

模型的整体前向传播可概括为：
\begin{equation}
    \hat{\mathbf{X}}, \mathbf{A}, \tilde{\mathbf{X}} = \mathrm{AGF\text{-}ADNet}(\mathbf{X}, \mathbf{M})
\end{equation}
其中 $\hat{\mathbf{X}} \in \mathbb{R}^{B \times S \times N}$ 为重构输出（即模型对输入信号的"预测"），$\mathbf{A} \in \mathbb{R}^{N \times N}$ 为学习到的邻接矩阵（描述变量间的关系），$\tilde{\mathbf{X}}$ 为插补中间结果（融合模块的输出，用于填补缺失数据）。

% ============================================================
\subsection{模块详解}

% ------------------------------------------------------------
\subsubsection{数据预处理模块（TEPDatasetBuilder）}

工业数据在输入模型之前需要经过仔细的预处理。本模块负责从 CSV 文件中加载\textbf{田纳西-伊斯曼过程}（Tennessee Eastman Process, TEP）数据集并进行标准化。TEP 是化工过程控制领域最经典的基准数据集之一，包含 41 个测量变量，涵盖温度、压力、流量等物理量。

\paragraph{为什么需要标准化？}

不同传感器的测量范围差异很大（例如温度可能在 $100$$\sim$$200\,{}^{\circ}\mathrm{C}$，而压力可能在 $1$$\sim$$10\,\mathrm{MPa}$）。如果直接使用原始数据，量纲较大的变量会主导模型的学习过程。Z-score 标准化通过减去均值、除以标准差，将每个变量映射到均值为 0、标准差为 1 的分布上：
\begin{equation}
    D_{t,n}^{\text{scaled}} = \frac{D_{t,n} - \mu_n}{\sigma_n}
\end{equation}
其中 $\mu_n, \sigma_n$ 分别为第 $n$ 个变量在训练集（前 1500 个样本）上的均值和标准差。

\paragraph{为什么要用训练集的统计量？}

这一点很重要：我们\emph{仅使用训练集}的均值和标准差来标准化\emph{所有}数据（包括测试集）。原因在于：在真实场景中，测试数据是未来到来的数据，我们无法提前知道其统计特性。使用训练集统计量可以模拟真实的在线部署场景，避免"信息泄露"。

\paragraph{缺失值处理。}

工业数据中经常存在缺失值（传感器故障、通信中断等）。模型通过构建缺失值掩码来处理：
\begin{equation}
    M_{t,n} = \mathbf{1}[D_{t,n} \text{ 非缺失}]
\end{equation}
其中 $\mathbf{1}[\cdot]$ 为指示函数，值为 1 或 0。缺失位置的数据被填充为零：$D_{t,n} \leftarrow D_{t,n} \cdot M_{t,n}$。

\paragraph{滑动窗口构造。}

深度学习模型通常需要固定长度的输入。为此，我们使用\textbf{滑动窗口}（Sliding Window）方法将长时间序列切分为等长片段。给定窗口长度 $S$ 和步长 $\delta$（默认 $\delta=1$），对于时间步 $t$，构造回看窗口：
\begin{equation}
    \mathbf{X}^{(t)} =
    \begin{cases}
        [\underbrace{\mathbf{d}_0, \ldots, \mathbf{d}_0}_{S-t-1}, \mathbf{d}_0, \ldots, \mathbf{d}_t] & \text{若 } t < S \\[6pt]
        [\mathbf{d}_{t-S+1}, \ldots, \mathbf{d}_t] & \text{若 } t \geq S
    \end{cases}
\end{equation}
其中 $\mathbf{d}_t \in \mathbb{R}^N$ 为第 $t$ 步的特征向量。当 $t < S$ 时（即序列开头，历史数据不足以填满窗口），通过复制首样本进行前端填充（padding），这是时间序列处理中的常用技巧。

% ------------------------------------------------------------
\subsubsection{自适应图学习模块（GraphLearner）}

\paragraph{直觉与动机。}

在一个化工厂中，41 个传感器并不是独立工作的——温度变化会影响压力，流量改变会引起液位变化。这些变量之间的关系可以用\textbf{图}（Graph）来表示：每个变量是一个\textbf{节点}（Node），如果两个变量之间存在依赖关系，就在它们之间连一条\textbf{边}（Edge），边的权重表示依赖的强度。

传统方法通常需要领域专家手动定义变量之间的关系图。GraphLearner 的创新之处在于：它让模型\textbf{自动学习}这个图结构，无需任何先验知识。

\paragraph{数学描述。}

GraphLearner 的输出为归一化邻接矩阵 $\mathbf{A} \in \mathbb{R}^{N \times N}$，其中 $A_{ij}$ 表示变量 $j$ 对变量 $i$ 的影响权重。

具体地，定义两组可学习的\textbf{节点嵌入}（Node Embedding）$\mathbf{E}_1, \mathbf{E}_2 \in \mathbb{R}^{N \times d_e}$（$d_e=16$）。嵌入是一种低维向量表示，每个节点被映射为一个 16 维向量，其含义由模型在训练过程中自动学习。

设放大因子 $\alpha=3.0$，邻接矩阵的计算分为三步：
\begin{enumerate}
    \item \textbf{非线性变换}：通过 $\tanh$ 函数将嵌入值压缩到 $[-1, 1]$ 范围，乘以 $\alpha$ 可以使分布更加"尖锐"，增强区分度：
    \begin{align}
        \mathbf{M}_1 &= \tanh(\alpha \, \mathbf{E}_1) \\
        \mathbf{M}_2 &= \tanh(\alpha \, \mathbf{E}_2)
    \end{align}

    \item \textbf{相似度计算}：通过矩阵乘法计算每对节点之间的"相似度"，并使用 ReLU 函数过滤掉负值（只保留正相关的依赖关系）：
    \begin{equation}
        \mathbf{A}^{\text{raw}} = \mathrm{ReLU}\!\left(\mathbf{M}_1 \, \mathbf{M}_2^\top\right)
    \end{equation}
    \textbf{解读}：$\mathbf{M}_1 \mathbf{M}_2^\top$ 的第 $(i,j)$ 个元素是节点 $i$ 的嵌入（来自 $\mathbf{E}_1$）与节点 $j$ 的嵌入（来自 $\mathbf{E}_2$）之间的内积。内积越大，说明两个节点越"相似"，它们之间的边权重就越大。使用两组不同的嵌入（$\mathbf{E}_1$ 和 $\mathbf{E}_2$）使得关系可以是非对称的，即"变量 $A$ 影响变量 $B$"不一定意味着"变量 $B$ 对变量 $A$ 有同等影响"。

    \item \textbf{行归一化}：为使邻接矩阵具有归一化的权重（类似概率分布），进行逐行归一化：
    \begin{equation}
        A_{ij} = \frac{A_{ij}^{\text{raw}}}{\sum_{k=1}^{N} A_{ik}^{\text{raw}} + \epsilon}
    \end{equation}
    其中 $\epsilon = 10^{-8}$ 为防止除零的数值稳定常数。
\end{enumerate}

\paragraph{为什么不使用 softmax 归一化？}

读者可能会问：归一化为什么不用更常见的 $\mathrm{softmax}$ 函数？原因在于 softmax 会使\emph{所有}权重都大于零（即图中所有节点两两相连），即使某对变量之间实际上没有关系。这使得后续通过 $L_1$ 正则化促进图稀疏性（即去除不必要的边）变得无效。使用 ReLU + 行归一化的组合，允许权重为零，使得正则化可以真正"剪掉"无关的边。

% ------------------------------------------------------------
\subsubsection{图卷积层（GCNLayer）}

\paragraph{直觉与动机。}

有了描述变量间关系的图 $\mathbf{A}$ 之后，我们需要一种方法来利用这些关系。图卷积网络（Graph Convolutional Network, GCN）的思想是：\textbf{每个节点的新表示是其邻居节点表示的加权平均}。

类比：想象一个教室里的学生（节点），如果学生 A 和学生 B 关系密切（边权重大），那么 A 的"综合表现"应该受到 B 的较大影响。GCN 就是在做类似的"消息传递"。

\paragraph{数学描述。}

输入为 $\mathbf{X} \in \mathbb{R}^{B \times S \times N \times F_{\text{in}}}$（$F_{\text{in}}=1$ 表示每个变量只有一个特征维度），输出为 $\mathbf{H} \in \mathbb{R}^{B \times S \times N \times F_{\text{out}}}$（$F_{\text{out}}=1$）。

图卷积的计算过程为：
\begin{equation}
    \mathbf{H} = \mathbf{A} \, (\mathbf{X} \, \mathbf{W} + \mathbf{b})
\end{equation}
其中 $\mathbf{W} \in \mathbb{R}^{F_{\text{in}} \times F_{\text{out}}}$ 为可学习权重矩阵，$\mathbf{b}$ 为偏置项。

展开来看每个元素的计算，使用 Einstein 求和约定：
\begin{equation}
    H_{b,s,n,f} = \sum_{m=1}^{N} A_{n,m} \cdot (\mathbf{X} \mathbf{W} + \mathbf{b})_{b,s,m,f}
\end{equation}

\textbf{解读}：对于节点 $n$（即第 $n$ 个传感器），其输出是所有节点经过线性变换后的加权和，权重就是邻接矩阵 $A_{n,m}$。如果节点 $m$ 和节点 $n$ 之间的边权重 $A_{n,m}$ 很大，那么节点 $m$ 的信息就会对节点 $n$ 的表示产生较大影响。这样，每个变量都能"看到"与它相关的其他变量的信息。

% ------------------------------------------------------------
\subsubsection{多尺度时序卷积网络（MultiScaleTCN）}

\paragraph{直觉与动机。}

工业过程中的异常模式可能发生在不同的时间尺度上：有些异常表现为瞬间的尖峰（短时尺度），有些则是缓慢的漂移趋势（长时尺度）。为了同时捕捉这些不同尺度的模式，我们使用多个不同大小的卷积核并行处理信号——核越大，感受野（Receptive Field）越大，能捕获越长期的模式。

\paragraph{什么是深度可分离卷积？}

为了减少参数量和计算开销，MultiScaleTCN 使用\textbf{深度可分离卷积}（Depthwise Separable Convolution）。普通卷积会同时在时间维度和通道维度上做卷积，参数量为 $O(N^2 \cdot k)$；深度可分离卷积将它拆分为两步：先对每个通道独立做时间卷积（$O(N \cdot k)$），然后通过 $1 \times 1$ 卷积混合通道信息。在本模型中只使用了第一步（分组卷积，$\text{groups}=N$），即每个变量独立地做时序卷积。

\paragraph{数学描述。}

输入张量 $\mathbf{X} \in \mathbb{R}^{B \times N \times S}$，输出张量 $\mathbf{Y} \in \mathbb{R}^{B \times N \times S}$。

设卷积核集合 $\mathcal{K} = \{3, 5, 9\}$，分别对应 3 步、5 步、9 步的时间窗口。对于每个核大小 $k_j$：
\begin{equation}
    \mathbf{Y}^{(j)} = \mathrm{DepthwiseConv1D}_{k_j}\!\left(\mathrm{Pad}(\mathbf{X}, k_j)\right), \quad j = 1, 2, 3
\end{equation}
其中 $\mathrm{Pad}(\cdot, k)$ 在本模型中使用对称填充（非因果模式，即同时利用过去和未来的信息）以保持序列长度不变。

对于核大小为 $k$ 的深度可分离卷积，第 $n$ 个变量在时间步 $s$ 的输出为：
\begin{equation}
    y_n^{(k)}(s) = \sum_{i=0}^{k-1} w_i^{(n,k)} \, x_n(s - \lfloor k/2 \rfloor + i) + b_n^{(k)}
\end{equation}
这本质上是一个滑动加权平均，其中权重 $w_i^{(n,k)}$ 是通过训练学习到的。

各尺度输出通过可学习权重向量 $\mathbf{w} \in \mathbb{R}^{|\mathcal{K}|}$ 融合：
\begin{equation}
    \mathbf{Y} = \sum_{j=1}^{|\mathcal{K}|} w_j \, \mathbf{Y}^{(j)} + b_{\text{fusion}}
\end{equation}
实现上，将三个尺度的输出堆叠为 $\mathbf{Y}^{\text{stack}} \in \mathbb{R}^{B \times N \times S \times |\mathcal{K}|}$，然后通过线性层（$|\mathcal{K}| \to 1$）对最后一维进行加权求和。模型会自动学习在不同场景下该赋予每个尺度多大的权重。

% ------------------------------------------------------------
\subsubsection{频域插补模块（FrequencyImputer）}

\paragraph{直觉与动机。}

到目前为止，所有的处理都是在\textbf{时域}（Time Domain）中进行的——直接对信号随时间变化的值进行操作。然而，信号还有另一种等价的表示方式：\textbf{频域}（Frequency Domain）。频域将信号看作多个不同频率的正弦波的叠加。

\textbf{为什么引入频域？}频域分析有几个独特的优势：
\begin{itemize}
    \item 周期性模式在频域中非常明显：一个频率 $f$ 的振荡在频谱中就是一个峰值；
    \item 噪声通常是高频成分，异常可能表现为特定频率的异常增强或消失；
    \item 时域中复杂的卷积运算在频域中变成简单的乘法，有助于全局信息的捕获。
\end{itemize}

\paragraph{离散傅里叶变换回顾。}

\textbf{傅里叶变换}（Fourier Transform）是将信号从时域转换到频域的数学工具。对于离散信号 $\{x_n(0), \ldots, x_n(S-1)\}$，离散傅里叶变换（DFT）定义为：
\begin{equation}
    X_n(f) = \sum_{s=0}^{S-1} x_n(s) \, e^{-j 2\pi f s / S}, \quad f = 0, 1, \ldots, \lfloor S/2 \rfloor
\end{equation}
其中 $j = \sqrt{-1}$ 为虚数单位。每个 $X_n(f)$ 是一个\textbf{复数}，可以分解为实部和虚部：$X_n(f) = \mathrm{Re}(X_n(f)) + j \, \mathrm{Im}(X_n(f))$，或者等价地，分解为\textbf{幅度}（Magnitude，$|X_n(f)|$）和\textbf{相位}（Phase，$\angle X_n(f)$）。幅度表示该频率成分的能量强度，相位表示其起始偏移。

逆变换将频域信号恢复为时域：
\begin{equation}
    x_n(s) = \frac{1}{S} \sum_{f=0}^{S-1} X_n(f) \, e^{j 2\pi f s / S}
\end{equation}

在实际实现中，使用快速傅里叶变换（FFT）算法，复杂度为 $O(S \log S)$。

\paragraph{频率注意力机制。}

并非所有频率成分对重构都同等重要。FrequencyImputer 使用\textbf{注意力机制}（Attention Mechanism）来自动学习每个频率成分的重要性，类似于人类注意力会选择性地关注重要信息。

输入 $\mathbf{X} \in \mathbb{R}^{B \times S \times N}$，处理过程如下：

\textbf{步骤 1：时域到频域。} 将输入转置为 $(B, N, S)$ 并执行 FFT：
\begin{equation}
    \mathbf{X}_f = \mathrm{FFT}(\mathbf{X}) \in \mathbb{C}^{B \times N \times F}, \quad F = \lfloor S/2 \rfloor + 1
\end{equation}

\textbf{步骤 2：构造特征。} 将复数频谱的实部与虚部拼接为实值特征向量：
\begin{equation}
    \mathbf{Z} = [\mathrm{Re}(\mathbf{X}_f); \, \mathrm{Im}(\mathbf{X}_f)] \in \mathbb{R}^{B \times N \times 2F}
\end{equation}

\textbf{步骤 3：计算注意力权重。} 通过一个两层的神经网络，为每个频率分量计算一个 $[0,1]$ 之间的注意力分数：
\begin{equation}
    \boldsymbol{\alpha} = \sigma\!\left(\mathbf{W}_2^{\text{att}} \, \mathrm{ReLU}\!\left(\mathbf{W}_1^{\text{att}} \, \mathbf{Z} + \mathbf{b}_1^{\text{att}}\right) + \mathbf{b}_2^{\text{att}}\right) \in [0,1]^{B \times N \times F}
\end{equation}
其中 $\sigma(\cdot)$ 为 Sigmoid 函数（将任意实数映射到 $[0,1]$），$\mathbf{W}_1^{\text{att}} \in \mathbb{R}^{128 \times 2F}$，$\mathbf{W}_2^{\text{att}} \in \mathbb{R}^{F \times 128}$。

\textbf{步骤 4：频率增强与加权。} 同时，通过另一个两层网络生成增强后的频率表示：
\begin{equation}
    \hat{\mathbf{Z}} = \mathbf{W}_2^{\text{enh}} \, \mathrm{ReLU}\!\left(\mathbf{W}_1^{\text{enh}} \, \mathbf{Z} + \mathbf{b}_1^{\text{enh}}\right) + \mathbf{b}_2^{\text{enh}} \in \mathbb{R}^{B \times N \times 2F}
\end{equation}
将增强后的实部和虚部分别与注意力权重相乘（逐元素，$\odot$ 表示 Hadamard 积）：
\begin{align}
    \hat{R}_f &= \hat{Z}_{[\cdot, \cdot, :F]} \odot \boldsymbol{\alpha} \\
    \hat{I}_f &= \hat{Z}_{[\cdot, \cdot, F:]} \odot \boldsymbol{\alpha}
\end{align}
注意力权重接近 1 的频率成分被保留，接近 0 的被抑制。

\textbf{步骤 5：频域回到时域。} 重构复数频谱并执行逆傅里叶变换：
\begin{equation}
    \mathbf{X}^{\text{freq}} = \mathrm{IFFT}\!\left(\hat{R}_f + j \, \hat{I}_f\right) \in \mathbb{R}^{B \times N \times S}
\end{equation}
最终转置回 $(B, S, N)$ 维度。

% ------------------------------------------------------------
\subsubsection{门控融合模块（GatedFusion）}

\paragraph{直觉与动机。}

我们现在有了两个分支的输出：时域分支 $\mathbf{H}^{\text{time}}$（捕获变量间的空间依赖和时序模式）和频域分支 $\mathbf{H}^{\text{freq}}$（捕获频率特征）。如何将它们结合起来？

最简单的方法是直接相加或取平均，但这假设两个分支在所有情况下都同等重要——实际上，某些时间步可能时域信息更有价值（如检测瞬时尖峰），另一些时间步频域信息更重要（如检测周期性异常）。

\textbf{门控融合}使用一个可学习的"门"（Gate）来动态决定每个时间步、每个变量应该更多地依赖哪个分支。门的值 $z \in [0,1]$ 控制混合比例：$z$ 接近 1 时偏向时域，$z$ 接近 0 时偏向频域。

\paragraph{数学描述。}

输入为 $\mathbf{H}^{\text{time}}, \mathbf{H}^{\text{freq}} \in \mathbb{R}^{B \times S \times N}$，输出为 $\mathbf{H}^{\text{fused}} \in \mathbb{R}^{B \times S \times N}$。

首先将两个分支沿变量维度拼接（转置至通道优先格式）：
\begin{equation}
    \mathbf{C} = [\mathbf{H}^{\text{time}}; \, \mathbf{H}^{\text{freq}}] \in \mathbb{R}^{B \times 2N \times S}
\end{equation}
通过时序卷积门控网络计算门控信号：
\begin{equation}
    \mathbf{Z} = \sigma\!\left(\mathrm{Conv1D}_{1 \times 1}\!\left(\mathrm{ReLU}\!\left(\mathrm{Conv1D}_{3}\!\left(\mathbf{C}\right)\right)\right)\right) \in [0,1]^{B \times N \times S}
\end{equation}
其中第一层卷积（核大小 3）将通道数从 $2N$ 压缩至 64，同时利用大小为 3 的核捕获局部时序上下文；第二层 $1\times 1$ 卷积将通道数恢复为 $N$。最终的融合输出为：
\begin{equation}
    \mathbf{H}^{\text{fused}} = \mathrm{LayerNorm}\!\left(\mathbf{Z} \odot \mathbf{H}^{\text{time}} + (1 - \mathbf{Z}) \odot \mathbf{H}^{\text{freq}}\right)
\end{equation}

\textbf{解读}：$\mathrm{LayerNorm}$（层归一化）的作用是稳定训练过程——它将输出归一化至均值为 0、方差为 1 的分布，防止网络中某些层的输出过大或过小导致梯度爆炸或消失。

% ------------------------------------------------------------
\subsubsection{Transformer 重构器}

\paragraph{直觉与动机。}

经过前面的模块，我们已经得到了融合了时空和频率信息的中间表示。最后一步是利用 \textbf{Transformer} 对整个时间序列进行全局建模并输出最终的重构信号。

Transformer 最初被提出用于自然语言处理（NLP），其核心是\textbf{自注意力机制}（Self-Attention）：序列中每个位置都可以"关注"（Attend to）所有其他位置，从而捕获长距离的依赖关系。与卷积的局部视野不同，自注意力天然具有全局感受野。

\paragraph{数学描述。}

输入为插补后的序列 $\mathbf{X}^{\text{filled}} \in \mathbb{R}^{B \times S \times N}$，输出为重构信号 $\hat{\mathbf{X}} \in \mathbb{R}^{B \times S \times N}$。

\textbf{步骤 1：输入投影与位置编码。}

首先通过线性投影将每个时间步的 $N$ 维特征映射到 $d=64$ 维的模型空间：
\begin{equation}
    \mathbf{Z}_0 = \mathbf{X}^{\text{filled}} \, \mathbf{W}^{\text{in}} + \mathbf{b}^{\text{in}} + \mathbf{P} \in \mathbb{R}^{B \times S \times d}
\end{equation}
其中 $\mathbf{W}^{\text{in}} \in \mathbb{R}^{N \times d}$ 为输入投影矩阵，$\mathbf{P} \in \mathbb{R}^{1 \times S \times d}$ 为可学习的位置编码。

\textbf{为什么需要位置编码？}自注意力机制是\textbf{排列不变的}——它对输入序列的顺序不敏感。但时间序列的顺序显然很重要（先发生的事件和后发生的事件意义不同）。位置编码通过为每个位置添加一个唯一的向量，让模型能够区分序列中不同时间步的信息。

\textbf{步骤 2：多头自注意力与前馈网络。}

Transformer 编码器由 $L=2$ 层堆叠而成，每层包含：
\begin{itemize}
    \item \textbf{多头自注意力}（Multi-Head Attention, $h=4$ 头）：将输入拆分为 4 个子空间，在每个子空间中独立计算注意力，然后拼接结果。多头机制让模型能够同时关注不同类型的依赖关系。
    \item \textbf{前馈网络}（Feed-Forward Network, FFN）：两层全连接网络（隐藏维度 128），对每个位置独立进行非线性变换。
\end{itemize}

每层的计算采用\textbf{残差连接}（Residual Connection）+ 层归一化的结构：
\begin{align}
    \mathbf{Z}_l' &= \mathrm{LayerNorm}\!\left(\mathbf{Z}_{l-1} + \mathrm{MultiHeadAttn}(\mathbf{Z}_{l-1})\right) \\
    \mathbf{Z}_l &= \mathrm{LayerNorm}\!\left(\mathbf{Z}_l' + \mathrm{FFN}(\mathbf{Z}_l')\right), \quad l = 1, 2
\end{align}

\textbf{残差连接}的作用是让梯度可以"跳过"某些层直接反向传播，有效缓解深层网络的训练困难（梯度消失问题）。

\textbf{步骤 3：输出投影。}

最终通过线性投影将 $d$ 维特征恢复至 $N$ 维变量空间：
\begin{equation}
    \hat{\mathbf{X}} = \mathbf{Z}_L \, \mathbf{W}^{\text{out}} + \mathbf{b}^{\text{out}} \in \mathbb{R}^{B \times S \times N}
\end{equation}
这就是模型最终的重构输出。

% ============================================================
\subsection{数据流}

AGF-ADNet 中数据的完整流动过程如下所述。

\subsubsection{输入预处理}

\begin{enumerate}
    \item 原始 TEP 数据经 Z-score 标准化后，按滑动窗口划分为 $\mathbf{X} \in \mathbb{R}^{B \times S \times N}$；
    \item 训练阶段引入\textbf{随机掩码增强}：以概率 $p=0.1$ 对已观测位置进行额外掩码，得到输入掩码 $\mathbf{M}^{\text{input}} = \mathbf{M} \odot \mathbf{R}$，其中 $R_{b,s,n} \sim \mathrm{Bernoulli}(1-p)$（即每个位置有 $1-p=0.9$ 的概率保留，$p=0.1$ 的概率被额外遮盖）；
    \item 输入数据按掩码置零：$\mathbf{X}^{\text{input}} = \mathbf{X} \odot \mathbf{M}^{\text{input}}$。
\end{enumerate}

\subsubsection{特征变换与分支处理}

\begin{enumerate}
    \item \textbf{图学习}：GraphLearner 生成邻接矩阵 $\mathbf{A} \in \mathbb{R}^{N \times N}$；
    \item \textbf{时域分支}：
    \begin{itemize}
        \item GCN：$\mathbf{X} \xrightarrow{\text{unsqueeze}} \mathbb{R}^{B \times S \times N \times 1} \xrightarrow{\text{GCN}(\mathbf{A})} \mathbb{R}^{B \times S \times N \times 1} \xrightarrow{\text{squeeze}} \mathbf{H}^{\text{gcn}} \in \mathbb{R}^{B \times S \times N}$，随后经 LayerNorm 归一化；
        \item TCN：$\mathbf{X} \xrightarrow{\text{permute}} \mathbb{R}^{B \times N \times S} \xrightarrow{\text{MultiScaleTCN}} \mathbb{R}^{B \times N \times S} \xrightarrow{\text{permute}} \mathbf{H}^{\text{tcn}} \in \mathbb{R}^{B \times S \times N}$；
        \item 残差叠加：$\mathbf{H}^{\text{time}} = \mathbf{H}^{\text{gcn}} + \mathbf{H}^{\text{tcn}}$；
    \end{itemize}
    \item \textbf{频域分支}：$\mathbf{X} \xrightarrow{\text{FrequencyImputer}} \mathbf{H}^{\text{freq}} \in \mathbb{R}^{B \times S \times N}$，经 LayerNorm 归一化。
\end{enumerate}

\subsubsection{融合与重构}

\begin{enumerate}
    \item \textbf{门控融合}：$\mathbf{H}^{\text{fused}} = \mathrm{GatedFusion}(\mathbf{H}^{\text{time}}, \mathbf{H}^{\text{freq}})$；
    \item \textbf{缺失值插补}：$\mathbf{X}^{\text{filled}} = \mathbf{X} \odot \mathbf{M} + \mathbf{H}^{\text{fused}} \odot (1 - \mathbf{M})$。这一步的直觉是：对于有观测值的位置，仍使用原始数据；对于缺失的位置，用模型的预测来填充；
    \item \textbf{Transformer 重构}：$\hat{\mathbf{X}} = \mathrm{Transformer}(\mathbf{X}^{\text{filled}})$。
\end{enumerate}

\subsubsection{异常评分与检测}

对于每个测试样本 $t$，异常分数定义为掩码加权的均方重构误差：
\begin{equation}
    \mathrm{Score}(t) = \frac{\sum_{s,n} \left(\hat{X}_{s,n}^{(t)} - X_{s,n}^{(t)}\right)^2 \cdot M_{s,n}^{(t)}}{\sum_{s,n} M_{s,n}^{(t)}}
\end{equation}

\textbf{解读}：分子是所有已观测位置上重构误差的平方和，分母是已观测位置的总数。本质上这就是一个\textbf{均方误差}（MSE），只是忽略了缺失数据的位置。

阈值由训练集上的异常分数统计量确定（基于\textbf{三倍标准差规则}）：
\begin{equation}
    \tau = \mu_{\text{train}} + 3 \, \sigma_{\text{train}}
\end{equation}
若 $\mathrm{Score}(t) > \tau$，则判定样本 $t$ 为异常。

\textbf{为什么使用三倍标准差？}在统计学中，对于正态分布的数据，约 99.7\% 的数据落在均值的三个标准差范围内（即"三西格玛规则"）。因此，超过 $\mu + 3\sigma$ 的分数极有可能不属于正常分布，这为我们提供了一个合理且有统计学依据的异常检测阈值。

% ============================================================
\subsection{训练目标与损失函数}

AGF-ADNet 的训练损失由三部分组成，每一部分都有其特定的作用。

\subsubsection{时域重构损失}

这是最基本的损失项，衡量模型重构信号与真实信号之间的差异。仅在已观测位置上计算均方误差：
\begin{equation}
    \mathcal{L}_{\text{time}} = \frac{\sum_{b,s,n} \left(\hat{X}_{b,s,n} - X_{b,s,n}\right)^2 \cdot M_{b,s,n}}{\sum_{b,s,n} M_{b,s,n} + \epsilon}
\end{equation}
这是一个标准的 MSE 损失，只是通过掩码 $\mathbf{M}$ 过滤掉了缺失数据的位置（因为我们无法评价模型在没有真实值的位置上的表现）。

\subsubsection{频域重构损失}

仅在时域中比较可能不够全面。频域损失强制模型在频谱上也保持与真实信号的一致性。对观测比例超过 50\% 的有效样本，比较重构信号与真实信号的\textbf{幅度谱}差异：
\begin{equation}
    \mathcal{L}_{\text{freq}} = \frac{1}{|\mathcal{B}_{\text{valid}}|} \sum_{b \in \mathcal{B}_{\text{valid}}} \frac{1}{NF} \sum_{n,f} \left(|\hat{X}_{b,n}(f)| - |X_{b,n}(f)|\right)^2
\end{equation}
其中 $\mathcal{B}_{\text{valid}} = \{b : \bar{M}_b > 0.5\}$ 为有效样本集合（观测率超过 50\% 的样本），$\bar{M}_b$ 为第 $b$ 个样本的平均观测率。

\textbf{为什么需要频域损失？}时域损失倾向于让模型学习到一个"平滑"的重构（因为 MSE 对尖锐变化的惩罚较大），而频域损失可以鼓励模型保留信号中的高频细节和周期性结构。两者互补，使重构质量更全面。

\subsubsection{图稀疏性正则化}

通过 $L_1$ 范数促进邻接矩阵的稀疏性：
\begin{equation}
    \mathcal{L}_{\text{sparse}} = \frac{1}{N^2} \sum_{i,j} |A_{ij}|
\end{equation}

\textbf{为什么需要稀疏性？}如果不加约束，图学习模块可能会学到一个"全连接"的图（所有变量两两相连），这不仅计算开销大，而且会引入噪声信息。$L_1$ 正则化鼓励大部分边权重趋向于零，只保留真正有意义的依赖关系。这与 Lasso 回归中的思想完全一致。

\subsubsection{总损失函数}

\begin{equation}
    \mathcal{L} = \mathcal{L}_{\text{time}} + \lambda_f \, \mathcal{L}_{\text{freq}} + \lambda_s \, \mathcal{L}_{\text{sparse}}
\end{equation}
其中 $\lambda_f = 0.1$ 和 $\lambda_s = 0.01$ 分别为频域损失和稀疏性正则化的权重系数。这些超参数的选择反映了各项的相对重要性：时域重构损失权重最大（系数为 1），因为直接的信号重构是最核心的目标；频域损失作为辅助约束，权重较小；图稀疏性正则化权重最小，因为它只是一个"软约束"。

\subsubsection{自监督训练策略}

AGF-ADNet 采用\textbf{自监督随机掩码策略}进行训练，这一策略的灵感来源于 NLP 领域的\textbf{掩码语言模型}（Masked Language Model，如 BERT）。核心思想是：随机遮盖一部分已知数据，训练模型恢复被遮盖的部分。

在每个训练批次中：
\begin{enumerate}
    \item \textbf{生成随机掩码}：$\mathbf{R} \sim \mathrm{Bernoulli}(0.9)^{B \times S \times N}$，即每个元素有 90\% 的概率为 1（保留），10\% 的概率为 0（被遮盖）；
    \item \textbf{构造输入掩码}：$\mathbf{M}^{\text{input}} = \mathbf{M} \odot \mathbf{R}$。这意味着原本就缺失的位置仍然缺失，而原本已观测的位置中有 10\% 被额外遮盖；
    \item \textbf{模型接收的输入}：被遮盖后的数据 $(\mathbf{X} \odot \mathbf{M}^{\text{input}}, \mathbf{M}^{\text{input}})$；
    \item \textbf{损失计算的关键}：损失在\textbf{所有原始已观测位置} $\mathbf{M}$ 上计算（而非仅在 $\mathbf{M}^{\text{input}}$ 上）。这意味着模型不仅要重构它"看到"的数据，还要恢复被随机遮盖掉的数据。这迫使模型学习数据的内在结构和规律，而不仅仅是简单地"复制"输入。
\end{enumerate}

\subsubsection{优化配置}

\begin{itemize}
    \item \textbf{优化器}：Adam（自适应矩估计），初始学习率 $\eta = 10^{-3}$。Adam 是目前深度学习中最常用的优化器之一，它结合了动量（Momentum）和自适应学习率的优点；
    \item \textbf{学习率调度}：ReduceLROnPlateau（当验证损失在 10 个 epoch 内不再下降时，将学习率乘以 0.5）。这种策略在训练初期使用较大的学习率快速收敛，后期逐步降低以精细调整；
    \item \textbf{梯度裁剪}：$\|\nabla\|_{\max} = 1.0$。当梯度的范数超过 1.0 时，按比例缩小。这防止了梯度爆炸（Gradient Explosion），尤其在处理时间序列和 Transformer 模型时非常重要；
    \item \textbf{训练轮次}：100 个 epoch，批量大小 32。一个 epoch 意味着模型遍历一次完整的训练集。
\end{itemize}

\end{document}
